% !PRINTED 313
\weberpage{313}

For the application we assemble the seven triples, from which one can read off the quadruples as the complement to the full system. In the first column stands the symbol ($\alpha$) to which the triple belongs:

\[
\begin{array}{c||ccc|}
1\,0\,0 & 0\,1\,0 & 0\,0\,1 & 0\,1\,1 \\
0\,1\,0 & 1\,0\,0 & 0\,0\,1 & 1\,0\,1 \\
0\,0\,1 & 1\,0\,0 & 0\,1\,0 & 1\,1\,0 \\
0\,1\,1 & 1\,0\,0 & 0\,1\,1 & 1\,1\,1 \\
1\,0\,1 & 0\,1\,0 & 1\,0\,1 & 1\,1\,1 \\
1\,1\,0 & 0\,0\,1 & 1\,1\,0 & 1\,1\,1 \\
1\,1\,1 & 0\,1\,1 & 1\,0\,1 & 1\,1\,0 \\
\end{array}
\]

If we denote the symbols ($\alpha$), in the order in which they stand in the first column of this table, by $1, 2, 3, 4, 5, 6, 7$, then we can represent the table for the triples and quadruples more simply thus:

\[
\begin{array}{c|ccc|cccc}
1 & 2 & 3 & 4 & 1 & 5 & 6 & 7 \\
2 & 1 & 3 & 5 & 2 & 4 & 6 & 7 \\
3 & 1 & 2 & 6 & 3 & 4 & 5 & 7 \\
4 & 1 & 4 & 7 & 2 & 3 & 5 & 6 \\
5 & 2 & 5 & 7 & 1 & 3 & 4 & 6 \\
6 & 3 & 6 & 7 & 1 & 2 & 4 & 5 \\
7 & 4 & 5 & 6 & 1 & 2 & 3 & 7 \\
\end{array}
\]

\begin{center}
\S.\,82.

\gesperrt{Application to equations of degree $8$.}
\end{center}

In the formulae derived above we now put for the variables $X_{z_1,z_2,z_3}$ the roots of an equation of degree $8$, of which we assume that its \gesperrt{Galois} group $P = TQ$, in the fixed domain of rationality $\Omega$, is contained in the linear congruence group $R = SQ$, so that $T$ is a divisor of $S$. All functions that admit the permutations of this group are rational.

To these belongs first of all the sum of the $X$:

\weq{1}{\sum^{z} X_{z_1, z_2, z_3} = B,}
%CONT>>
